\section{Methods}
\label{sec:method}

\subsection{Architectures}

\paragraph{Students.}
Our student family is a post-activation basic-block ResNet. A $3\times3$ stem
produces the first stage's channels and keeps the $32\times32$ resolution; three
stages then use 3/5/3 or 2/3/3 basic blocks, with the first block of stages 2 and
3 applying stride 2. Each block computes
$\mathrm{ReLU}(F(x)+S(x))$, where $F$ is conv--BN--ReLU--conv--BN, the last BN is
zero-initialized~\citep{goyal2017} so the residual branch starts at zero and the block starts close
to a shortcut, and $S$ is an identity or a
$1\times1$ conv--BN shortcut when the shape changes. A global average pool and a
linear layer produce the logits. The two selected presets are
\texttt{r3\_64\_3-5-3} (64/128/256 channels, 3/5/3 blocks, 4.92M parameters) and
\texttt{r3\_80\_2-3-3} (80/160/256, 2/3/3, 4.88M). A 300-epoch sweep also
evaluated \texttt{r3\_96\_2-2-2} (4.90M), \texttt{r4\_64\_3-3-2-1} (4.89M), and
\texttt{r3\_40\_8-8-8} (4.72M).

\paragraph{Design choices.}
Inside the 5M budget a stage costs roughly $B_i C_i^2$ parameters, so depth and
width trade off directly. We keep $3\times3$ convolutions ($5\times5$ costs
$2.8\times$ as much, and two stacked $3\times3$ convs already cover a
$5\times5$ receptive field), $1\times1$ projection shortcuts only where the
resolution halves (a $3\times3$ shortcut would spend budget on the pass-through
path), no max-pool in the stem so stage 1 runs at full resolution, and global
average pooling over the final $8\times8$ map. Wider stages run faster on a GPU
and give each block more features, but every block is expensive; deeper,
narrower nets are slower and become thin early (40 channels at full resolution
for \texttt{r3\_40}). The middle ground \texttt{r3\_80} (8 blocks,
80--256 channels) beat both the deeper \texttt{r3\_64} (11 blocks) and the
wider \texttt{r3\_96} (6 blocks) at 300 epochs; that mid-resolution capacity
matters most is a hypothesis from this small sweep.

\paragraph{Teachers.}
We use two WRN-28-10~\citep{zagoruyko2016wide} teachers (36.5M parameters each, seeds 0 and 1) and one
PyramidNet-110-270~\citep{han2017pyramidal} teacher with ShakeDrop~\citep{yamada2019shakedrop} (28.5M, seed 2). Teachers are trained
on the 45,000 training images for 400 epochs with the same augmentation recipe and a different seed each; the WRNs use SAM, the PyramidNet does not;
teacher size is not constrained by the 5M budget because only the student is
submitted. The teacher pool reaches 98.14\% holdout accuracy (WRN seeds 0+1); the
individual teachers reach 98.06, 98.12, and 97.64.

\subsection{Training Recipe}

All students use SGD with Nesterov momentum 0.9, batch size 128, base learning
rate 0.1, weight decay $5\times10^{-4}$ (not applied to BN affine parameters or
biases), a 5-epoch linear warmup followed by cosine decay to zero, and bf16
autocast. The default augmentation is random crop with 4-pixel padding plus
horizontal flip, TrivialAugment, Cutout 16, and per-batch mixup ($\alpha=0.2$) or
CutMix ($\alpha=1.0$) with probability 1.0. Mixing is disabled for the last 15
epochs of 300-epoch runs, 30 epochs of 600-epoch runs, 50 epochs of 1,200-epoch
runs, and 20 epochs of teacher runs. Label smoothing~\citep{szegedy2016rethinking} is 0.1 when training from hard labels and is disabled when
KD is on.

SAM with $\rho=0.05$ costs two forward/backward passes per step; the reported SAM
runs use it for the whole schedule. EMA with decay 0.999 is maintained alongside
the raw weights, and both are evaluated every epoch. The 300-epoch ablation
schedule is the shipping schedule for the augmentation comparison; the main
students use 600 or 1,200 epochs.

\subsection{Knowledge Distillation}

For ordinary KD, each teacher sees the same augmented and mixed batch as the
student (consistent teaching~\citep{beyer2022patient}), and its softened distribution is computed at temperature $T=4$. The
teacher distributions are averaged, and the student minimizes
\begin{equation}
\mathcal{L} = 0.9\,T^2\,\mathrm{KL}\!\left(\bar{p}_T \,\|\, q_T\right)
+ 0.1\,\mathrm{CE}(q, y),
\end{equation}
where $\bar{p}_T$ is the averaged teacher distribution and $q_T$ the student
distribution at the same temperature. The CE term uses the mixup/CutMix soft
targets. The two-teacher (WRN seed 0 + PyramidNet) and three-teacher runs use this loss with
\texttt{r3\_64}; the selected \texttt{r3\_80} runs use the same loss with all
three teachers.

We also implement logit-standardized KD (LS-KD)~\citep{sun2024logit}. Each teacher's logits are
z-scored per example before softening, the ensemble is averaged in that space, and
the student's logits are standardized identically. The loss is
$0.1\,\mathrm{CE} + 9\,T^2\,\mathrm{KL}$ with $T=2$. We evaluate LS-KD both as a
60-epoch continuation and as a fresh 300-epoch run.

\subsection{Augmentation Policies}

We compare five alternatives to the TrivialAugment control, all at the same
150-epoch pilot schedule (no SAM, no KD):
\begin{itemize}
\item \textbf{RandAugment}~\citep{cubuk2020randaugment}, two operations per image.
\item \textbf{AutoAugment-CIFAR}~\citep{cubuk2019autoaugment}, the learned CIFAR policy.
\item \textbf{Geometry only}: mild random affine (translation, scale, rotation).
\item \textbf{Extra affine+color}: stronger affine plus color jitter.
\item \textbf{AugMix}: the torchvision implementation with severity 3 and mixture
width 3, one view and no JSD term. This is the image-only variant; it does not add
the consistency loss.
\end{itemize}
The winner is then evaluated at 300 epochs with SAM, replicated with a second
seed, and finally tested on top of three-teacher KD. This schedule matters: the
150-epoch pilot overstates the augmentation effect because the 300-epoch control
also improves substantially.
